\scanpage{346}
\mspage{\deleted{\unsure{12|10}}11}

déduite de $\unsure{\Upsilon|Y}$ en lui rajoutant un objet $\unsure{o_0|\sigma_0}$, relié à $a_0$ par deux flèches distinctes $\lambda_0,\mu_0$ telles que
\[
\beta_0\lambda_0=\beta_0\mu_0\qquad\text{et}\qquad\gamma_0\lambda_0=\gamma_0\mu_0,
\]
ces deux flèches composées étant, avec les quatre flèches $\beta_0,\gamma_0,\lambda_0,\mu_0$ marquées sur le diagramme, les seules (six) flèches non identiques de $Z$. $\Upsilon$ \deleted{\unsure{s’identifie|identifiée}} est immergée dans $Z$ \deleted{\unsure{comme un|pour un}} un fermé, de complémentaire la catégorie ponctuelle réduite à $o_0$. (N.B. L’immersion fermée $\Upsilon\longrightarrow Z$ n’est pas immersion de Dwyer \inserted{(\unsure{comme} $Z$ est le plus petit ouvert de $Z$ contenant $\Upsilon$)} i.e. n’est pas coasphérique, car la catégorie $o\backslash\Upsilon$, cofibrée sur $\Upsilon$ avec fibres discrètes, de cardinal $2$ en $a_0$ et $1$ en $b_0$ et $c_0$, est la catégorie
\diagram{\begin{tikzcd}[ampersand replacement=\&,row sep=small,column sep=large]
a'_0\arrow[r]\arrow[dr] \& \widetilde b_0\\
a''_0\arrow[ur]\arrow[r] \& \widetilde c_0
\end{tikzcd}}
, \deleted{\unsure{qui}} \inserted{i.e.}
\diagram{\begin{tikzcd}[ampersand replacement=\&,row sep=tiny,column sep=tiny]
{} \& \bullet\\
\bullet\arrow[ur]\arrow[dr] \& {} \& \bullet\arrow[ul]\arrow[dl]\\
{} \& \bullet
\end{tikzcd}}
\inserted{, qui} a le type d’homotopie du cercle). Ni l’immersion ouverte $\{o_0\}\unsure{\text{---}|\longrightarrow}Z$, car $\{o_0\}/a_0$ est la cat. discrète de card. $2$.

\deleted{\unsure{Z}}L’inclusion \deleted{\unsure{de}} $i:Z_0\text{---}Z$ de
\diagram{\begin{tikzcd}[ampersand replacement=\&,column sep=small]
Z_0={} \& o_0\arrow[r,bend left=20,"\lambda_0"]\arrow[r,bend right=20,"\mu_0"'] \& a_0
\end{tikzcd}}
dans $Z$ \unsure{est} \deleted{\unsure{une} \unsure{équivalence}} homotopisme\textsuperscript{*} (car il existe une rétraction évidente $r$ de $Z$ sur $Z_0$\textsuperscript{*}, et une flèche $i\unsure{r}\longrightarrow\mathrm{id}_Z$), donc $Z$ est \unsure{homotope} à \unsure{$Z$|$Z_0$}, qui a le type d’homotopie d’un cercle, donc n’est pas $1$\inserted{$\unsure{2}$}-connexe,
% slanted
\margin{faux, il n’y a pas de rétraction du tout ! En fait, $Z$ est la \uline{suspension} de
\diagram{\begin{tikzcd}[ampersand replacement=\&,row sep=tiny,column sep=small]
{}\arrow[r,bend left=25]\arrow[r,bend right=25] \& {}
\end{tikzcd}}
, qui \unsure{a le|est de} type d’homotopie \unsure{(topologique)} du cercle, \unsure{et} \unsure{le type} d’homotopie de \unsure{la} $2$-sphère.}

% archive sheet number: 346
